\documentclass[10pt,a4paper]{amsart}
\usepackage[margin=19mm]{geometry}
\usepackage{amsmath,amssymb,mathtools}
\usepackage[scr=rsfs,cal=euler]{mathalfa}
\usepackage{xfrac}
\usepackage[unicode,hidelinks,bookmarks=false]{hyperref}
\newcommand{\ie}{i.e.\ }
\newcommand{\cf}{cf.\ }
\newcommand{\resp}{respectively\ }
\newcommand{\Subsection}{\subsection}
\providecommand{\nobreakdash}{\nobreak\hbox{-}}
\setlength{\emergencystretch}{2em}
\multlinegap0pt
\usepackage{amssymb}
\usepackage{enumerate}
\usepackage{tensor}
\usepackage{breqn}
\usepackage{xcolor}
\usepackage{mathtools}
\usepackage{braket}
\usepackage{bm}
\usepackage{algorithm}\usepackage{algpseudocode}
\newtheorem{lemma}{Lemma}
\title{Invariant Theory, Magic State Distillation, and Bounds on Classical Codes}
\author{Amolak Ratan Kalra, Shiroman Prakash}
\date{Source: arXiv:2501.10163v3 (CC BY 4.0)}
\begin{document}
\maketitle

\noindent\emph{Source and licence.} Invariant Theory, Magic State Distillation, and Bounds on Classical Codes. Version arXiv:2501.10163v3 (CC BY 4.0), distributed by arXiv under the Creative Commons Attribution 4.0 International licence, http://creativecommons.org/licenses/by/4.0/, as recorded in the arXiv licence metadata for this exact version and shown on the abstract page https://arxiv.org/abs/2501.10163v3. Full licence notice: \texttt{licenses/arxiv-2501.10163-CC-BY-4.0.txt}.\par\smallskip

\noindent\textit{Editorial scope.} Counted unit: the article's lemma logical-operator-weight-n-lemma (source0.tex lines 400--403). The article's complete original proof of it is delivered with it (source0.tex lines 404--418, one \texttt{proof} environment of 15 lines). No mathematics is written, repaired or re-derived: the statement and the proof are the article's own lines, in the article's own notation, and no retained line is substituted or re-flowed.  Retained uncounted as the source-local context the proof needs: lines 371--373.  Nothing was omitted from any retained span, so the package carries no omission record.  No retained line contains a citation, so the wrapper carries no bibliography and no bibliography entry.  No retained line contains a figure environment, an image-inclusion command or drawing code, so no image is delivered and the package declares no asset dependency.  Editorial formatting only, in addition: an AMS \texttt{amsart} preamble that restates the article's own declarations and macros used by the retained lines, namely the environments \texttt{lemma} on separate counters, together with the standard packages those lines need.  The retained environments print in this document as Lemma 1 in order of appearance, whereas the article prints the same items with the numbering of its own full document.  The complete native source of the article as distributed in the arXiv e-print for this exact version is bundled unchanged as \texttt{sources/171727/source0.tex}.
\begin{equation}
    B(x,y) = \frac{1}{2^{n-k}} A(x+3y, x-y), \label{eq:classical-macwilliams}
\end{equation}

\begin{lemma}
\label{logical-operator-weight-n-lemma}
Every $[[n,1]]$ $M_3$-code possesses a logical operator of weight $n$.
\end{lemma}

\begin{proof}
Let $A(x,y) = \sum A_j x^{n-j}y^j$ be the simple weight enumerator of the classical code corresponding to the stabilizer group $\mathcal{S}$. By Rall's rule, all stabilizers have even weight, so $A_j = 0$ for all odd $j$.

The weight enumerator of the dual code, $B(x,y) = \sum B_j x^{n-j}y^j$, determines the weights of operators in the normalizer $N(\mathcal{S})$. Using the MacWilliams identity (Eq. \eqref{eq:classical-macwilliams}), the number of weight-$n$ operators in the dual is given by the coefficient of $y^n$:
\begin{equation}
    B_n = B(0,1) = \frac{1}{2^{n-1}} A(3, -1).
\end{equation}
Expanding the polynomial $A(3, -1)$:
\begin{equation}
    B_n = \frac{1}{2^{n-1}} \sum_{j=0}^n A_j 3^{n-j} (-1)^j.
\end{equation}
Since $A_j = 0$ for all odd $j$, the term $(-1)^j$ is always equal to $+1$ for all non-zero $A_j$. Furthermore, since $A_0=1$ and $A_j \geq 0$, the sum is strictly positive. Thus, $B_n > 0$.

Since $n$ must be odd for any $[[n,1]]$ $M_3$-code, this weight-$n$ operator cannot be a stabilizer (as stabilizers must have even weight). Therefore, it must be a logical operator.
\end{proof}

\end{document}
